\section{征服不是箭头：从草原联盟到长江水网}

\subsection{一二〇六年以后，称号先要变成供给}

一二〇六年诸部推举铁木真为成吉思汗，草原上的新称号并没有立刻变成一部可以在万里之外征税的机器。《元朝秘史》把联盟、分配与亲属关系放在叙事中央；《元史》本纪则从后来王朝的尺度回望。两种文本都提示了一个较为稳固的事实：新的政治中心必须把原先分散的随从、俘获物、牧地和使者编进可持续的关系。它首先面对的不是“中国”这一预定目标，而是马群的冬夏移动、部众的分配、通向绿洲和金朝边境的商路，以及谁能够向军队提供消息。

草原军队的速度容易使人忽略其停顿。马匹需要草场，向导需要熟悉水源，战利品与投附者需要重新分配；一支部队穿过沙漠边缘和河谷时，能否得到粮食、修理器具和翻译，常常比地图上两点的距离更重要。蒙古材料中的“降”可能指首领改变关系，也可能掩盖城镇、部落和家庭在不同时间作出的不同选择。因而，称号统一了命令的名义，却没有取消地方社会对道路、市场与水源的控制。后来的征服能展开，正在于蒙古统治者不断寻找能把这些地方条件接入军政网络的中介。

这也解释了为什么早期战争中畏兀儿、契丹、女真、汉人与中亚商人不能只作为“被征服者”出现。有人提供文书技术，有人熟悉城镇贸易，有人掌握边地语言，有人以投附换取新的保护。这样的合作并非天然忠诚，更不是单向的文化吸收；它发生在战争、家族安全、旧债务和新机会相互重叠的场景里。正史较容易记录首领的来归，较难记录没有姓名的运输者、牧人和逃离者。叙事只能确认前者的年序，却应把后者留下的沉默当作征服成本的一部分。

\subsection{西夏：河西走廊不是一条空走廊}

一二〇九年蒙古进攻西夏，战场的关键并不只在中兴府城下。黄河灌溉的农田、河套边缘的牧地、河西走廊的绿洲和通向西域的道路，使西夏既能成为蒙古南下前的障碍，也能成为连接多种语言、税源和宗教机构的区域。蒙古军围攻中兴府并引河水灌城的记载，保留了攻守的一项动作；堤决与撤围同样表明水并非征服者可以任意调遣的武器。河水、堤岸、季节、城内仓储与军队停留时间共同限制了围城。

此后的和议、贡赋和联姻，在王朝叙事中常被写成胜负已经分明的仪式。对中兴府、河西绿洲和寺院网络而言，事情更像一串未完成的交接：谁来征粮，旧官如何留用，商旅是否仍能通行，边地军队向谁请饷。西夏本国连续国史没有保存下来，汉文正史多从敌手或后继者的角度书写；黑水城等地的西夏文、汉文、藏文和回鹘文文书虽能让局部的法律、书写和宗教活动浮现，却不能填满每一座城、每一个村落的经历。正因此，不能把一二二七年的覆亡写成区域社会突然空白。

成吉思汗去世与西夏覆亡同年，后世很容易将二者并成一幕完成征服的戏剧结尾。较可靠的读法是，蒙古的军事与继承网络在这一刻都面临重新安排：河西的路线与税源变得重要，西夏遗留的官员、工匠、僧侣和商人也必须在新的保护关系中寻找位置。对后来的元朝而言，这些地区不是遥远附属物；它们是与中亚、吐蕃和华北相接的交通与宗教空间。对西夏居民而言，国号的终结也不自动说明语言、信仰或土地关系同步终结。

\subsection{野狐岭到中都：华北城防怎样变成接收难题}

一二一一年蒙古大举攻金，野狐岭及其后的北方战事使金朝原有的防线承受连续压力。史书喜欢列举攻取的州县，仿佛城名按次序落在一张清单上；但居庸关、紫荆关、燕山山口和河北平原之间的道路，决定军队能够如何分兵、粮车能够走到哪里、难民与守军又会怎样堵塞同一条通道。金廷既要保住中都，也要维持南方与东北的联系；蒙古军则要把投附的军人、地方首领和后勤资源组织为不随一次战役消失的力量。

一二一五年中都陷落后，最急迫的问题从破城转为帑藏、仓储、工匠和人口如何被接收。《元史》记有籍中都帑藏的行动，这类记录可以证实征服者试图把战利品转化为可管理的资源；它不能告诉我们每一项物资是否安全入库，更不能把被登记的人群写成没有选择的数字。城市中有人逃离，有人留守，有人借旧关系向新军请求保护；城外乡村还要面对田地、征发、市场和治安的变化。接收不是战争结束后的附录，而是决定征服能否持续的下一场考验。

华北各地的投降亦有不同含义。地方将领选择归附，可能是因为补给已断、旧主难以援救，也可能是为了保住部众和城镇；蒙古统帅接受他们，同样要承担监督、赏赐和防叛的成本。把所有投附者写成见风使舵，或把他们都写成忠于新秩序，都将当时的资源约束抹掉。金、蒙古和宋方材料分别突出自己的功绩、背叛或灾难，军数与伤亡尤其不宜照录。可以追踪的是道路、城池和仓储的控制怎样改变；不可轻率断言的是每一个居民的立场与损失。

\subsection{蔡州的短暂联盟，及其后的河南}

一二三四年蒙古与南宋共同攻破蔡州，金朝终结。这是一项最容易被后来结局误读的合作：双方都希望削弱金，却并没有共享同一幅北方秩序的蓝图。蔡州周围的围城、粮道与冬季行军，使协作带有明确而短暂的军事目的；城破之后，河南故地、边界与赋源立即成为新的竞争对象。宋方与蒙古方关于战功和责任的叙述不能互相抵消，反而显示联盟本身就充满解释权的争夺。

对河南居民，金亡没有自动带来一个可预期的新行政框架。旧军队、地方豪强、流民和新进入的军政人员要围绕城镇、道路和粮食重新定位。国家的更替先以征发、治安、田地和市场的具体改变抵达家庭，而非以史书中的断代标题抵达。墓志、碑刻与后来的地方材料偶尔能看见迁居、任官或重建，却偏向有资源留下文字的人；不能由少数仕宦家族的继续活动，推论整个河南已经恢复。

这场联盟的制度得失正在这里。它暂时解决了攻金所需的军事协调，却把如何分配战后资源的问题推迟下来。对于南宋，北方边界的变动加重了淮河与四川方向的防务顾虑；对于蒙古，接收华北意味着必须用更稳定的文书、税粮和地方中介维持军队。之后的宋蒙冲突并非蔡州一役的机械后果，却在同一批土地、渡口与粮道上获得了更尖锐的条件。

\subsection{大理、吐蕃与西南：山路上的帝国并不整齐}

一二五三年蒙古征服大理，常在中原年表中被压成一行。这一行掩盖了横断山、河谷、城镇与多种地方政治之间的距离。军队进入西南不只需要击败一个王朝，也要面对山地道路的季节性、马匹与粮食的补给、地方首领的协商以及不同宗教网络的存在。大理的接收改变了南宋西南的战略环境，但“已下”并不表示山地村落、坝区城镇与通向东南亚的道路已被同一种行政方式覆盖。

蒙古统治者在西南和吐蕃边地往往借助地方首领、僧侣与既有机构维持关系。这不应被概括为单纯的宽容或单纯的征服策略：不同中介掌握土地、劳力、仪式权威和翻译能力，中央也要用封授、驿传、军队和礼物不断维系这些关系。藏文、八思巴字、汉文与地方碑刻提供的只是片段，说明多语书写在行政与宗教之间流动；它们不能将任何一种语言等同于固定的政治忠诚。

一二五八年至一二五九年蒙哥发动攻宋计划并在四川前线去世，更能看出西南空间的限制。钓鱼城附近的山城、水路与补给，使大规模部署无法仅凭中央意志推进。宋、蒙两方对战事和死因的叙述各有立场，不能据此虚构完整的战场细节；可以确认的是统帅的去世迫使原有战争安排中断，并把继承与资源调度问题带回更广阔的帝国。山地前线并非边缘背景，它直接改变了草原、华北和江南政治的时间。

\subsection{兄弟争位与山东叛乱：谁控制了运粮线}

一二六〇年忽必烈即位后，同阿里不哥的争位延续至一二六四年。两都、草原与华北的不同资源网络使继承不能只在宫廷中解决：支持者需要粮食、马匹、驿路和能够传递命令的官员，城市和地方将领也要判断何一方能兑现赏赐与保护。把冲突写成两个性格的较量，便会看不见后勤和地域如何决定政治选择。战争结束后，胜者仍须重新接入原来被动员过的官员、军队与地方财政。

李璮在山东的叛乱尤其揭示了接收与自治的紧张。山东连接华北、海运与南宋前线，地方武装和粮食运输不只是中央可随时调用的资源。叛乱遭镇压，说明忽必烈集团试图限制能够独立掌握军政与财政的地方中介；但镇压本身也要依赖同样的道路、仓储与合作网络。制度上得到的是对军事地方化的警惕，代价则是中枢必须更直接地承担监督、供给与任官的负担。

此时的“汉地”也不能被当成同质区域。河北、山东、山西与河南经历战乱、投附和恢复的节奏不同，社会上留下的军户、工匠、旧官与流动人口也不同。后来制度文本中的分类，记录的是统治者试图整理这种复杂性的方式，不能反向证明每个人已经被稳定地放进一个类别。要理解忽必烈的统治为何转向大都、行省与更强的文书网络，必须先看到继承战争和地方叛乱留下的这一难题。

\subsection{襄樊到崖山：长江水网拒绝一战定局}

一二六八年开始的襄阳、樊城围攻，延续到一二七三年。两城之所以重要，不是因为它们在地图上恰好位于宋元之间，而是汉水、荆襄道路、船只、粮运和江南财政在这里相接。围城军需要阻断补给，也需要维持自己的补给；守军需要依靠水路与周边网络，城内居民则要在食物、疾病、劳役和撤离机会中承受长期压力。宋、元叙事对攻守功绩各有安排，攻城器械与援军的细节不能脱离其文类照录。

襄樊陷落后，元军南下并非一条直入临安的道路。四川、两淮、长江中游和东南沿海具有不同的地形、船只和地方组织。一二七六年临安降服可以确立为政治节点，却不能替代各处接收：有的地方官改事新朝，有的宗室和将领转向海上，有的市镇要先处理仓粮、户籍与治安。南宋遗民诗文所记的丧失感是真实的文化证词，却不能让士人的声音吞没船户、农民、工匠和军人的不同处境。

一二七九年崖山海战后，南宋政权终结。后世常以幼主与海面收束叙事，但对广东、福建和江南港口而言，战争留下的是船只、贸易、沿海防务和人口流动的持续重组。港口遗存、碑刻与地方文书能够让个别航线和社群浮现，不能量化所有海上损失。征服在此得到的，是通往江南税粮和海域网络的更大可能；它承担的代价，是必须把新接收地区的运输、法律和地方关系纳入一个本来就跨越草原与农耕区的帝国。
